%% ====================================================================
\section{Introduction and Motivation}
\label{sec:intro}
%% ====================================================================

\subsection{The Riemann Hypothesis}
\label{sec:rh}

The Riemann zeta function is defined for $\re(s) > 1$ by the absolutely convergent
series
\begin{equation}\label{eq:zeta-def}
  \zeta(s) = \sum_{n=1}^{\infty} n^{-s},
\end{equation}
and extended to the entire complex plane (except for a simple pole at $s=1$) by
analytic continuation. The functional equation
\begin{equation}\label{eq:func-eq}
  \zeta(s) = \chi(s)\,\zeta(1-s),
  \qquad
  \chi(s) = 2^s \pi^{s-1} \sin\!\bigl(\tfrac{\pi s}{2}\bigr)\,\Gamma(1-s),
\end{equation}
equivalently expressed using the completed zeta function
\begin{equation}\label{eq:xi-def}
  \xi(s) = \tfrac{1}{2}\,s(s-1)\,\pi^{-s/2}\,\Gamma\!\bigl(\tfrac{s}{2}\bigr)\,\zeta(s),
\end{equation}
which satisfies the symmetric form $\xi(s) = \xi(1-s)$, implies that $\zeta(s)$ vanishes
at $s = -2, -4, -6, \ldots$ (the \emph{trivial zeros}). The Riemann
Hypothesis~\cite{riemann1859} asserts:

\medskip
\noindent\textbf{RH.} \emph{All nontrivial zeros of $\zeta(s)$ have real part equal
to $\tfrac{1}{2}$.}
\medskip

\noindent
Equivalently, if $\zeta(\sigma + it) = 0$ with $0 < \sigma < 1$ and $t \in \R$, then
$\sigma = \tfrac{1}{2}$. This remains the most important open problem in mathematics,
with profound implications for the distribution of prime numbers, the error term in the
prime counting function, and numerous results across number theory, analysis, and
mathematical physics that are conditional on~RH.

\subsection{Computational Approaches to RH}
\label{sec:computational}

The dominant computational approach to~RH has been direct numerical verification:
computing $\zeta(\tfrac{1}{2} + it)$ along the critical line and certifying that all
zeros up to some height~$T$ lie on it. Odlyzko~\cite{odlyzko1989,odlyzko2001} computed
specific high zeros of~$\zeta$---around the $10^{20}$th and $10^{22}$nd zeros---and
large blocks of consecutive zeros near these heights, providing detailed statistical
data on zero spacing and pair correlation. Gourdon~\cite{gourdon2004} systematically
verified that all of the first $10^{13}$ nontrivial zeros lie on the critical line,
using the Odlyzko--Sch\"onhage algorithm based on the Riemann--Siegel formula. These
computations exploit highly optimized algorithms specific to the zeta function and
achieve extraordinary precision.

Our approach differs fundamentally. Rather than computing $\zeta$ directly, we
\emph{classically} evaluate the mathematical expressions (partial Dirichlet sums,
Loschmidt amplitudes, rate functions) that arise from quantum-mechanical systems
whose observable phase transitions correspond exactly to the zeros
of~$\zeta$.\footnote{To be precise: no quantum hardware or quantum simulation is
involved. The computation reduces to evaluating partial sums of Dirichlet series and
related functions on a classical computer. The quantum-mechanical language (Hamiltonians,
coherence functions, Loschmidt echoes) provides the \emph{framework} and
\emph{interpretation}, not the computational method. The value of this framing lies in
connecting the Riemann Hypothesis to the well-developed theory of dynamical quantum phase
transitions, not in any computational advantage over direct zeta evaluation.} This
provides:
\begin{enumerate}
  \item An independent computational pathway to zero detection, serving as a cross-check.
  \item Physical intuition about the structure of zeta zeros through the language of
        quantum phase transitions.
  \item A connection (though not a direct realization) to the Hilbert--P\'olya
        program\footnote{The Hilbert--P\'olya conjecture, traditionally attributed to
        c.~1914 (though the earliest written record is later), posits the existence of
        a self-adjoint operator whose eigenvalues are the imaginary parts of the
        nontrivial zeta zeros. If such an operator exists, the reality of its
        eigenvalues would immediately imply~RH.}, linking~RH to quantum dynamical
        phenomena rather than to a single spectral operator.
\end{enumerate}

\subsection{The Hilbert--P\'olya Conjecture and DQPTs}
\label{sec:hilbert-polya}

Berry and Keating~\cite{berrykeating1999} proposed that the Hilbert--P\'olya operator
might be related to the quantization of the classical Hamiltonian $H = xp$, connecting~RH
to quantum chaos.

The Wei et~al.\ construction~\cite{wei2025} takes a different but related path. Instead
of seeking a single operator whose spectrum encodes the zeros, they construct quantum
systems whose \emph{dynamical} phase transitions---non-analyticities in a rate function
as a function of real time---occur precisely at the zeta zeros. The key insight is that
the Loschmidt amplitude of these systems, when evaluated at inverse temperature~$\beta$,
is proportional to the Dirichlet eta function $\eta(\beta + it)$ (System~1) or the
Hardy $Z$-function (System~2). The~RH is then equivalent to the statement that DQPTs
occur only at $\beta = \tfrac{1}{2}$.

This reformulation shifts~RH from a statement about complex analysis to one about
quantum statistical mechanics, opening new avenues for both theoretical and numerical
investigation.

\subsection{Scope and Limitations of the DQPT Reformulation}
\label{sec:scope}

We emphasize an important conceptual point: the DQPT--RH equivalence established by
Wei et~al.\ is a \emph{reformulation} of the Riemann Hypothesis in the language of
quantum statistical mechanics, not a proof mechanism. The equivalence states that~RH
holds if and only if DQPTs are confined to $\beta = \tfrac{1}{2}$, but this
reformulation does not by itself provide any new leverage toward proving (or
disproving)~RH. The difficulty of the problem is preserved: it is merely re-expressed
in physical language.

In particular:
\begin{enumerate}
  \item The quantum systems (System~1 and System~2) are \emph{constructed} so that
        their partition functions and order parameters reproduce the Dirichlet eta
        function and Hardy $Z$-function. The physics does not generate~RH as a
        consequence of some independent physical principle; rather, the mathematics
        of~RH is \emph{encoded} into the physics by design.
  \item The numerical investigation described in this document can provide computational
        evidence consistent with (or inconsistent with)~RH, but it cannot prove~RH.
        This is a fundamental limitation shared by all numerical approaches.
  \item The value of the DQPT framework lies in (a)~providing physical intuition about
        the structure of zeta zeros, (b)~connecting~RH to the well-developed theory of
        quantum phase transitions, and (c)~potentially suggesting new analytical
        approaches inspired by the physics (see Section~\ref{sec:directions}).
\end{enumerate}
